%
% 0_Einleitung.tex -- Einleitung zum Kapitel Resultante
%
% (c) 2026 Jakob Gierer, OST Ostschweizer Fachhochschule
%
% !TEX root = ../../buch.tex
% !TEX encoding = UTF-8
%
\section{Einleitung\label{resultante:section:einleitung}}
\kopfrechts{Einleitung}

Für die Lösung von Integralen greifen Ingenieur:innen sowie Studierende oft auf Computeralgebrasysteme (CAS) zurück.
Damit CAS-fähige Rechner Integrale lösen können, werden effiziente Algorithmen benötigt.
Integranden werden zunächtst in einen polynomiellen Teil und rationalen Teil aufgeteilt.
Dieses Kapitel behandelt nur Methodem zur Bestimmung der Stammfunktion des rationalen Teils.
Dabei besteht der rationale Teil des Integrals aus echt gebrochenen Funktionen. 
Für die Berechnung von diesen Integralen, müssen die Linearfaktoren/Nullstellen des  des Integranden bestimmte werden.

In Abschnitt \ref{buch:ratint:section:rothsteintrager} wird der Satz von Rothstein-Trager beschrieben, der für die Bestimmung der Nullstellen in Kapitel \ref{chapter:intalg} den euklidischen Algorithmus nutzt.
In diesem Kapitel soll eine weitere Methode zur Berechnung von Nullstellen vorgestellt werden, die Resultante.
In verschiedenen Textbücheren wird die Resultante als
\begin{equation}
    \operatorname{Res}(f,g)=\det(\operatorname{Syl}(f,g))
    \label{resultante:equation:definition}
\end{equation}
definiert. Aber was ist \(\operatorname{Syl}\)? 
Und was ist die Resultante?

Was genau die Resultante ist und was für nützliche Eigenschafte sie besitzt, beschreibt dieses Kapitel.